\chapter{Utilizzo del robot reale UR5}
\label{chap:ur5_reale}

Dopo aver sviluppato e testato il sistema di controllo del robot in ambiente simulato con Unity, l'implementazione è stata estesa per consentire l'utilizzo del robot reale UR5.
Questa sezione descrive le modifiche apportate all'architettura e al flusso di lavoro per adattarsi alle specifiche del robot reale, garantendo un'integrazione efficace con il sistema di controllo basato su ROS~2.

\section{Polyscope e test preliminari}
Prima ancora di passare al robot reale, abbiamo fatto utilizzo di un ulteriore simulatore, ovvero Polyscope, che ci ha permesso di testare il sistema di controllo effettivo in un ambiente virtuale più vicino alla realtà, simulando il comportamento del robot UR5 e verificando la correttezza delle traiettorie pianificate e dei comandi inviati al robot.
Su Polyscope abbiamo utilizzato il driver UR5 ROS2 di default, ma un differente controller, ovvero il \texttt{forward\_position\_controller}, che ci ha permesso di utilizzare lo stesso metodo precedente per la trasmissione della traiettoria pianificata, ovvero la pubblicazione dei riferimenti punto per punto ad una certa frequenza.

\section{Test sul robot reale}
Dopo il processo precedente, il sistema è stato testato sul robot reale e le traiettorie pianificate sono state eseguite con successo, dimostrando l'efficacia del sistema di controllo sviluppato sebbene il metodo di pubblicazione dei riferimenti punto per punto ad una certa frequenza, anche elevata, non sia stata ottimale e presentava dei leggeri scatti durante l'esecuzione della traiettoria.
Motivo per il quale, è stato deciso di passare ad un differente controller ros2 che è quello di default del driver UR5, ovvero lo \texttt{scaled\_joint} e che utilizza un \textit{action server} e un \textit{action client} per la trasmissione della traiettoria pianificata, permettendo di inviare l'intera traiettoria in un'unica richiesta e ricevere una conferma di completamento al termine dell'esecuzione.

\subsection{Scelta del controller: \texttt{scaled\_joint\_trajectory\_controller}}

Il driver ROS~2 del robot UR mette a disposizione diversi controller. I due più rilevanti per questo lavoro sono il \texttt{forward\_position\_controller} e lo \texttt{scaled\_joint\_trajectory\_controller}.

Il \texttt{forward\_position\_controller} è un controller di tipo \emph{forward command}: riceve su un topic un vettore di posizioni articolari e lo inoltra direttamente all'interfaccia hardware.
Non esegue alcuna interpolazione, non applica limiti di velocità o accelerazione e non verifica il raggiungimento dell'obiettivo.
Il profilo di moto dipende interamente dal nodo che pubblica i comandi, che deve quindi inviare un nuovo setpoint a frequenza costante e senza interruzioni.
Un ritardo o un salto tra due setpoint successivi si traduce in un movimento brusco del robot.

Lo \texttt{scaled\_joint\_trajectory\_controller} è invece un controller di traiettoria: riceve, tramite l'action \texttt{FollowJointTrajectory},
una sequenza di punti (posizioni, velocità e istanti temporali) e si occupa di interpolarli e di eseguirli. Il termine \emph{scaled} indica che il tempo della traiettoria viene scalato in base al fattore di velocità
impostato sul teach pendant del robot: se l'operatore riduce la velocità o interviene una sospensione di sicurezza, la traiettoria rallenta
mantenendo la propria forma, senza accumulare errore di inseguimento.

\begin{table}[h]
\centering
\begin{tabular}{p{4cm}p{4.5cm}p{4.5cm}}
\toprule
 & \texttt{forward\_position} & \texttt{scaled\_joint\_trajectory} \\
\midrule
Ingresso & singolo setpoint (topic) & traiettoria completa (action) \\
Interpolazione & assente & presente \\
Temporizzazione & a carico del nodo & gestita dal controller \\
Tolleranze e esito & assenti & presenti (feedback e risultato) \\
Annullamento & non previsto & previsto (cancellazione del goal) \\
Scalatura di velocità & no & sì (teach pendant) \\
\bottomrule
\end{tabular}
\caption{Confronto tra i due controller.}
\end{table}

Si è scelto lo \texttt{scaled\_joint\_trajectory\_controller} per tre motivi.
In primo luogo, la traiettoria viene completamente calcolata \emph{offline}
a partire dalla traiettoria cartesiana e inviata come un unico goal: il controller
garantisce la temporizzazione, mentre con il \texttt{forward\_position\_controller}
la regolarità del moto dipenderebbe dalla periodicità di un nodo Python
non real-time. In secondo luogo, il controller fornisce feedback, verifica
delle tolleranze e possibilità di annullare il movimento, utili per
controllare l'esecuzione e fermare il robot in caso di anomalie. Infine, la
scalatura del tempo rispetta i limiti e gli arresti di sicurezza imposti
dal pendant, aspetto rilevante in un'applicazione in cui il robot opera
in prossimità dell'operatore. Infine, grazie all'interpolazione interna, il controller garantisce un movimento più fluido e privo di scatti, migliorando la qualità complessiva dell'esecuzione della traiettoria pianificata.